% реферат по теме 'АИС и АСУ'
\documentclass{SibFU-docs}

% доп. пакеты
\usepackage{graphicx}
\usepackage{float}
\usepackage{hyperref}
\usepackage{caption}
\usepackage{subcaption}
\usepackage{amsmath}
% предотвартим ошибку с \Bbbk перед загрузкой пакета amssymb
\let\Bbbk\relax
\usepackage{amssymb}
\usepackage{booktabs}
\usepackage{tabularx}
\usepackage{longtable}
\usepackage{enumitem}
\usepackage{multirow}

% гиперссылки
\hypersetup{
    colorlinks=true,
    linkcolor=black,
    citecolor=blue,
    urlcolor=blue
}

% библиография
\addbibresource{report.bib}

% титульный лист
\begin{document}
\setlist[itemize]{left=1.3cm, itemsep=1pt, topsep=2pt} % для выравнивания списков
\setlist[enumerate]{left=1.3cm, itemsep=1pt, topsep=2pt}
\renewcommand{\contentsname}{}
\mycovertitle
{Гуманитарный институт}
{Кафедра прикладной информатики в искусстве и интерактивном медиа}
{Технологии создания и обработки графических данных}
{ст. преподаватель И.~Р.~Нигматуллин}
{Логинова Ю.В., Захаров И.М., Пачковская А.В., Финочка С.А., \\Ткачева М.С., Мельников С.В., Черепанова А.А., Тимофеева А.А.}
 
\begin{center}
\bfseries РЕФЕРАТ
\end{center}
\noindent

Реферат по теме: «Технологии создания и обработки графических данных» содержит 31 страниц и 18 использованных источников.

Ключевые слова: РАСТРОВАЯ ГРАФИКА, ВЕКТОРНАЯ ГРАФИКА, ЦВЕТОВЫЕ МОДЕЛИ, RGB, CMYK, ГЛУБИНА ЦВЕТА, ФОРМАТЫ ФАЙЛОВ, JPEG, PNG, SVG, СЖАТИЕ ДАННЫХ, РЕТУШЬ, КОРРЕКЦИЯ ЦВЕТА, ФИЛЬТРЫ, МЕТАДАННЫЕ, ТЕГИ, СТРУКТУРА ПАПОК, РЕЗЕРВНОЕ КОПИРОВАНИЕ, КАТАЛОГИЗАЦИЯ, ЦИФРОВОЙ АРХИВ, БЕЗОПАСНОЕ ХРАНЕНИЕ.

В работе исследуются фундаментальные принципы, методы и технологии, лежащие в основе всего цикла работы с цифровыми изображениями. Рассмотрены математические и структурные основы растрового и векторного представления графики, теория и практика цветовых моделей, алгоритмы сжатия и форматы хранения данных. Проанализированы основные инструменты и приемы редактирования, ретуши и улучшения изображений. Систематизированы подходы к организации, каталогизации и долговременному архивному хранению графических коллекций в цифровой среде.

Основные выводы:
\begin{enumerate}
\item Растровая и векторная графика представляют собой принципиально разные математические модели, каждая из которых оптимальна для своих задач: растровая для фотореалистичных изображений, векторная для логотипов, чертежей и типографики.
\item Цветовая модель является системообразующим фактором и определяется областью применения: RGB для экранов, CMYK для полиграфии, что требует грамотного преобразования для сохранения цветового охвата.
\item Выбор формата файла и метода сжатия (с потерями или без) представляет собой компромисс между качеством изображения, размером файла и требованиями к его дальнейшему использованию и редактированию.
\item Современные графические редакторы реализуют единые фундаментальные принципы обработки (работа со слоями, выделениями, кривыми), владение которыми важнее знания конкретного интерфейса.
\item Цифровая обработка изображений является целенаправленным итеративным процессом коррекции дефектов и оптимизации визуальных характеристик для достижения конкретного коммуникативного результата.
\item Организация и хранение графических коллекций требуют системного подхода, сочетающего логическую структуру папок, работу с метаданными и строгое соблюдение правила резервного копирования 3 2 1 для гарантии сохранности данных.
\end{enumerate}

Рекомендации:
\begin{itemize}[label=-]
\item При создании нового графического актива первичным должен быть выбор между растровой и векторной моделью в зависимости от требуемого результата и масштабируемости.
\item Работать с изображением следует в самом широком цветовом пространстве (например, Adobe RGB или ProPhoto RGB), конвертируя в конечную модель (sRGB, CMYK) только на этапе финального экспорта.
\item Для хранения исходников и промежуточных этапов редактирования использовать форматы без потерь (PSD, TIFF, RAW), для публикации в вебе применять оптимизированные форматы с потерями (JPEG, WebP).
\item Осваивать графические редакторы на уровне понимания общих принципов и инструментов (слои, маски, корректирующие слои), а не только запоминания расположения кнопок.
\item Внедрять систему каталогизации на основе метаданных (ключевые слова IPTC) с момента импорта файлов, что в разы ускоряет их поиск в будущем.
\item Создавать и автоматизировать процесс резервного копирования графических архивов по правилу 3 2 1, рассматривая его как обязательную часть рабочего процесса.
\end{itemize}

\newpage

\thispagestyle{empty}
\begin{center}
\bfseries СОДЕРЖАНИЕ
\end{center}
\vspace{-6pt}
\tableofcontents
\clearpage

% фантомные боли для секций, чтобы не было нумераций кроме тела реферата
\section*{}
\vspace*{-2.5em}
\begin{center}\textbf{ВВЕДЕНИЕ}\end{center}
\phantomsection
\addcontentsline{toc}{section}{Введение}

Графические данные являются одним из фундаментальных типов информации в цифровую эпоху, пронизывая все сферы — от веб-дизайна и медиа до научной визуализации и цифрового искусства. Технологии их создания и обработки формируют основу для коммуникации, творчества и анализа.

Современное состояние проблемы. Несмотря на повсеместную доступность графических редакторов, существует существенный разрыв между умением пользоваться отдельными инструментами и глубоким пониманием базовых принципов работы с цифровым изображением. Это приводит к неоптимальному выбору форматов, потере качества, неэффективной организации материалов и, как следствие, к снижению продуктивности и риску утраты цифровых активов. Ключевой задачей становится системное овладение полным циклом работы с графикой — от выбора правильной математической модели и цветового пространства до профессиональной обработки и обеспечения долговременной сохранности данных.

Актуальность темы обусловлена тотальной визуализацией информации, необходимостью в профессиональных навыках создания и управления графическим контентом в любой цифровой профессии, а также потребностью в сохранении целостности и доступности графических коллекций как части цифрового наследия.

Новизна исследования заключается в комплексном рассмотрении полного технологического цикла работы с графическими данными — от их генезиса (растровые и векторные модели) через обработку до архивации, с акцентом на практическое понимание взаимосвязи между теорией (цветовые модели, сжатие) и прикладным результатом.

Цель работы: исследовать ключевые технологии, методы и принципы создания, обработки, хранения и управления графическими данными в цифровой среде.

Задачи исследования: раскрыть принципиальные различия между растровой и векторной графикой и определить области их эффективного применения; проанализировать структуру, принципы работы и сферы использования основных цветовых моделей; исследовать популярные форматы графических файлов, принципы сжатия изображений и критерии их выбора; рассмотреть основы редактирования графических данных в прикладных программах; изучить основные методы цифровой обработки изображений для коррекции и улучшения их качества; систематизировать принципы и методы организации, каталогизации и безопасного долговременного хранения графических коллекций.

Методы исследования: сравнительный анализ для противопоставления растрового и векторного подходов, форматов и цветовых моделей; структурный анализ для изучения архитектуры файлов и алгоритмов сжатия; практический анализ интерфейсов и инструментария основных графических редакторов; систематизация методов на основе принципов информационного менеджмента и архивного дела применительно к цифровым коллекциям.

\newpage
% вручную увеличить номер секции и записать в оглавление с номером,
% чтобы в теле не печатался автоматический номер, но в tableofcontents он остался
\refstepcounter{section}
\addcontentsline{toc}{section}{\protect\numberline{\thesection} Отличие векторной и растровой графики и области их применения}
\vspace*{-2.5em}
\begin{center}\textbf{\thesection\ ОТЛИЧИЕ ВЕКТОРНОЙ И РАСТРОВОЙ ГРАФИКИ И ОБЛАСТИ ИХ ПРИМЕНЕНИЯ}\end{center}
\vspace*{-1.5em}
\ManualSubsection{Принципиальная основа и ключевые характеристики растровой графики}

Растровая графика представляет собой фундаментальный подход к представлению цифрового изображения, основанный на использовании пиксельной матрицы. Основной принцип заключается в формировании изображения из множества дискретных элементов — пикселей (pixels), каждый из которых обладает строго определёнными координатами в прямоугольной сетке и значением цвета. Совокупность этих элементов, воспринимаемая визуально как единое целое, и составляет растровую картину. Важнейшей характеристикой, напрямую влияющей на детализацию и качество, является разрешение, измеряемое в количестве пикселей на единицу длины (чаще всего точек на дюйм, DPI). Как отмечается в источниках, именно разрешение и глубина цвета являются ключевыми параметрами для растровых файлов \cite{grafika_formats}.

Главным ограничением растрового подхода выступает проблема масштабируемости. Поскольку изображение жестко привязано к фиксированной сетке пикселей, его значительное увеличение приводит к визуальному проявлению этой дискретной структуры — эффекту пикселизации, когда контуры становятся ступенчатыми, а общая четкость снижается. При этом основным достоинством растровой графики остаётся её исключительная способность к передаче сложных, непрерывных тоновых и цветовых переходов, что делает её незаменимой для фотореалистичных изображений, цифровой фотографии и живописи. Типичными форматами для хранения растровых данных являются JPEG, PNG, GIF и TIFF, каждый из которых оптимизирован под определённые задачи сжатия и хранения \cite{rastr_vs_vector}.

\ManualSubsection{Принципиальная основа и ключевые характеристики векторной графики}

Векторная графика основана на принципиально ином, математическом способе описания изображения. Вместо фиксации цвета каждой точки, векторный подход определяет графические примитивы (точки, линии, кривые, многоугольники) с помощью математических формул, например, через координаты опорных точек и параметры соединяющих их кривых Безье. Данный метод обеспечивает так называемую разрешающую независимость, что является ключевым преимуществом векторных изображений \cite{vektor_osnovy}.

Бесконечная масштабируемость без потери качества — прямое следствие математической природы векторной графики. При изменении размера изображения пересчитываются только формулы, описывающие объекты, что гарантирует идеально чёткие контуры при любом масштабе. Кроме того, каждый объект в векторном файле остаётся независимым и легко редактируемым, что обеспечивает высокую гибкость на этапе дизайна. Объём файла при этом зависит не от физических размеров картинки, а от сложности сцены и количества составляющих её объектов. Наиболее распространёнными форматами векторной графики являются SVG, AI, EPS и PDF (который может содержать как векторные, так и растровые данные).

\ManualSubsection{Сравнительный анализ областей применения растровой и векторной графики}

Выбор между растровым и векторным подходом является стратегическим и целиком определяется конкретной прикладной задачей. Области применения растровой графики преимущественно связаны с задачами, требующими работы с непрерывными тонами и максимальной детализацией. Это цифровая и печатная фотография, создание текстур и фонов для веб-сайтов, цифровая живопись, ретушь фотографий, а также обработка сложных визуальных эффектов в кинематографе. Растровая графика служит основным инструментом там, где важна точная передача реального визуального мира.

Векторная графика, в свою очередь, доминирует в сферах, где на первый план выходят требования к точности, масштабируемости и чёткости контуров. К ним относится разработка фирменного стиля и логотипов, которые должны одинаково качественно воспроизводиться на носителях любого размера — от визитной карточки до рекламного щита. Вектор незаменим в полиграфическом и веб-дизайне для создания макетов, иконок, интерфейсных элементов и инфографики. Он также лежит в основе современных шрифтов (TrueType, OpenType) и активно используется в системах автоматизированного проектирования (CAD) для создания инженерных чертежей и схем \cite{rastr_vs_vector}.

\ManualSubsection{Взаимодействие и преобразование форматов}

На практике часто возникает необходимость во взаимном преобразовании графических форматов, однако процессы конвертации вектора в растр и обратно имеют фундаментально разную природу и результат. Растеризация, или преобразование векторного изображения в растровое, представляет собой технически простой и контролируемый процесс. Программа вычисляет цвет каждого пикселя итоговой растровой сетки на основе математического описания векторных объектов, сохраняя идеальное качество для заранее заданного выходного разрешения. Этот процесс является необратимым и используется на финальных стадиях подготовки макетов к публикации или печати.

Обратная операция — векторизация (трассировка) растрового изображения — является значительно более сложной и зачастую неточной. Алгоритмы трассировки пытаются анализировать массив пикселей, выявляя в нём контуры и одноцветные области для последующего описания векторными кривыми. Качество результата сильно зависит от сложности исходного растра: простые изображения с чёткими контурами (логотипы) могут быть векторизованы успешно, в то время как попытка трассировки фотографии приведёт к созданию чрезмерно сложного и нефункционального набора векторных объектов. Таким образом, векторизация является не универсальным методом преобразования, а специализированным инструментом для решения узкого круга задач, таких как восстановление логотипа из растрового прототипа.

\newpage
\refstepcounter{section}
\addcontentsline{toc}{section}{\protect\numberline{\thesection}Структура и принципы работы цветовых моделей}
\begin{center}\textbf{\thesection\ СТРУКТУРА И ПРИНЦИПЫ РАБОТЫ ЦВЕТОВЫХ МОДЕЛЕЙ}\end{center}
\vspace*{-1.5em}
\ManualSubsection{Сущность и классификация цветовых моделей}

Цветовая модель представляет собой математическую систему описания цвета, позволяющую представить его через числовые значения параметров или метод смешения базовых компонентов. Разложение цвета на составляющие необходимо для его точного воспроизведения, хранения и обработки в цифровых системах. В основе классификации цветовых моделей лежат различные системы цветности, каждая из которых соответствует определённому физическому принципу или способу восприятия. Аддитивные модели, такие как RGB, основаны на сложении световых излучений, где отсутствие света даёт чёрный цвет, а смешение всех компонентов в максимуме — белый. Субтрактивные модели, например CMYK, оперируют понятием поглощения (вычитания) световых волн красящими веществами, где базовым белым цветом является цвет незапечатанной поверхности. Отдельный класс составляют перцепционные модели, такие как HSL и LAB, которые проектируются с учётом особенностей человеческого цветовосприятия, стремясь обеспечить более интуитивную и психофизиологически адекватную работу с цветом \cite{colorimetry2020}.

\ManualSubsection{Эталонная модель CIE XYZ}

Международная комиссия по освещению (CIE) в 1931 году разработала фундаментальную цветовую модель CIE XYZ, которая с тех пор служит эталоном в области колориметрии. Данная модель является строго математическим представлением, основанным на экспериментальных данных о среднем человеческом цветовосприятии. Она оперирует тремя гипотетическими, нереализуемыми физически компонентами (X, Y, Z), которые охватывают весь спектр видимого излучения. Модель CIE XYZ стала отправной точкой и универсальным языком для сравнения всех прочих цветовых пространств, а также основой для систем калибровки и профилирования устройств ввода-вывода, таких как мониторы, принтеры и сканеры \cite{colorimetry2020}.

\ManualSubsection{Аддитивная модель RGB и её применение}

Модель RGB является доминирующей аддитивной системой, используемой в устройствах, излучающих свет. Её структура основана на трёх основных цветах: красном (Red), зелёном (Green) и синем (Blue). Каждый цвет в этой модели кодируется тройкой числовых значений, обычно в диапазоне от 0 до 255, где (0,0,0) соответствует чёрному цвету (отсутствие излучения), а (255,255,255) — белому (максимальная интенсивность всех трёх компонентов). Принцип работы RGB основан на сложении световых потоков различных длин волн, что напрямую соответствует физике работы жидкокристаллических и светодиодных экранов, проекторов и цифровых камер. Именно поэтому данная модель установлена по умолчанию для подготовки контента, предназначенного для отображения на экранах \cite{rgb2022}.

\ManualSubsection{Субтрактивная модель CMYK и её применение}

Модель CMYK является стандартом в полиграфии и относится к субтрактивным системам. Она описывает цвет через процентное содержание четырёх типографских красок: голубой (Cyan), пурпурной (Magenta), жёлтой (Yellow) и чёрной (Key). Чёрный компонент был введён в модель по практическим соображениям: смешение трёх основных цветовых красок в теории должно давать чёрный, но на практике приводит к получению грязно-коричневого оттенка; кроме того, использование отдельной чёрной краски снижает стоимость печати и нагрузку на бумагу. Принцип работы CMYK основан на селективном поглощении белого света, падающего на бумагу: краска вычитает (поглощает) определённые части спектра, а отражённый свет формирует воспринимаемый цвет. Перевод изображения из RGB в CMYK является критически важным этапом препресса, так как цветовой охват печати уже, чем у экрана \cite{cmyk2021}.

\ManualSubsection{Перцепционные модели LAB и HSL}

Перцепционные цветовые модели проектируются с целью максимального соответствия человеческому восприятию, отделяя информацию о яркости от информации о цветовом тоне и насыщенности. Модель LAB, также известная как CIELAB, является аппаратно-независимой и строится на основе эталонной модели CIE XYZ. Она состоит из компонента L, обозначающего светлоту (яркость), и двух цветоразностных компонентов: «a» (ось от зелёного к красному) и «b» (ось от синего к жёлтому). Ключевым преимуществом LAB является возможность коррекции яркости без влияния на цветовой тон, что делает её незаменимой в профессиональной ретуши и цветокоррекции. Модель HSL (Hue, Saturation, Lightness) представляет цвет в виде интуитивно понятного цилиндра, где тон (Hue) задаётся углом на цветовом круге, насыщенность (Saturation) — расстоянием от центра, а светлота (Lightness) — положением по вертикальной оси. Этот подход, используемый во многих графических редакторах для выбора цвета, позволяет пользователю оперировать субъективными категориями, такими как «сделать цвет более ярким» или «менее насыщенным», что упрощает процесс дизайна.

\newpage
\refstepcounter{section}
\addcontentsline{toc}{section}{\protect\numberline{\thesection}Популярные форматы графических файлов и способы сжатия изображений}
\vspace*{-1.5em}
\begin{center}\textbf{\thesection\ ПОПУЛЯРНЫЕ ФОРМАТЫ ГРАФИЧЕСКИХ ФАЙЛОВ И СПОСОБЫ СЖАТИЯ ИЗОБРАЖЕНИЙ}\end{center}
\vspace*{-1.5em}
\ManualSubsection{Методы сжатия графических данных}

Сжатие графических данных является критически важным процессом, направленным на уменьшение объёма файлов для экономии дискового пространства и ускорения их передачи по сетям. Все методы сжатия разделяются на две фундаментальные категории: сжатие без потерь (Lossless) и сжатие с потерями (Lossy). Алгоритмы без потерь, такие как RLE (Run-Length Encoding), LZW (Lempel–Ziv–Welch) и DEFLATE, обеспечивают полное восстановление исходных данных при распаковке. RLE эффективен для изображений с большими однородными областями, заменяя последовательности одинаковых пикселей парой «значение–количество повторений» \cite{rle-wiki}. Более сложный алгоритм DEFLATE, объединяющий LZ77 и кодирование Хаффмана, лежит в основе форматов PNG и ZIP. В противоположность этому, сжатие с потерями, наиболее ярко представленное стандартом JPEG, намеренно отбрасывает часть информации, малозаметную для человеческого восприятия, что позволяет достигать существенно больших коэффициентов сжатия для фотографических изображений \cite{jpeg-standard}.

\ManualSubsection{Артефакты сжатия с потерями}

Использование сжатия с потерями неизбежно приводит к появлению визуальных искажений, особенно при высоких степенях сжатия. Наиболее характерными артефактами являются блочность (проявление квадратной структуры 8x8 пикселей, используемой в дискретном косинусном преобразовании JPEG), общее размытие мелких деталей и текстур, а также возникновение муара на регулярных паттернах и цветовых ореолов (бахромы) вокруг контрастных границ объектов. Эти искажения являются прямым следствием процессов квантования и прореживания, в ходе которых алгоритм жертвует высокочастотными компонентами изображения для сокращения его размера.

\ManualSubsection{Векторные форматы файлов}

Векторные форматы описывают изображение с помощью математических формул, определяющих геометрические примитивы, что обеспечивает их принципиальную масштабируемость без потери качества. Открытый стандарт SVG (Scalable Vector Graphics), основанный на XML, широко используется в веб-разработке для логотипов, иконок и инфографики, поддерживая сжатие без потерь алгоритмом GZIP. Универсальный формат PDF (Portable Document Format) применяет гибридный подход, сохраняя векторные данные и текст без потерь, в то время как растровые элементы могут сжиматься с потерями (JPEG). Специализированные проприетарные форматы, такие как AI (Adobe Illustrator) и CDR (CorelDRAW), предназначены для хранения сложных рабочих проектов со всеми слоями и эффектами.

\ManualSubsection{Классические растровые форматы}

Среди классических растровых форматов доминирующее положение занимает JPEG, стандартизированный объединённой группой экспертов по фотографии (Joint Photographic Experts Group). Он использует сложный многоэтапный алгоритм, включающий преобразование цветового пространства, дискретное косинусное преобразование и квантование, что позволяет эффективно сжимать фотографии с приемлемыми визуальными потерями \cite{jpeg-standard}. Формат PNG (Portable Network Graphics) был создан как свободная альтернатива GIF и использует алгоритм DEFLATE для сжатия без потерь, поддерживая полноценную прозрачность (альфа-канал) и глубину цвета до 48 бит, что делает его идеальным для графики с чёткими границами. Устаревший формат GIF, ограниченный 256 цветами, до сих пор используется для простой анимации, а профессиональный формат TIFF (Tagged Image File Format) сохраняет максимальное качество и метаданные, поддерживая различные схемы сжатия, и широко применяется в полиграфии и архивации.

\ManualSubsection{Современные и специализированные форматы}

Эволюция технологий привела к появлению новых, более эффективных форматов. WebP от Google сочетает преимущества JPEG и PNG, предлагая как сжатие с потерями (с лучшей эффективностью, чем у JPEG), так и без потерь, а также поддержку анимации и прозрачности, что способствует ускорению загрузки веб-страниц \cite{webp-guide}. Формат HEIF/HEIC, основанный на видеокодеках, обеспечивает лучшее качество при том же размере файла, что и JPEG, и активно продвигается в экосистеме Apple. Наиболее перспективным считается формат AVIF, построенный на открытом кодеке AV1, который демонстрирует наивысшую эффективность сжатия среди современных решений. Для профессиональной работы используются «сырые» форматы RAW, содержащие необработанные данные с матрицы камеры, и проприетарные редакторские форматы, такие как PSD (Adobe Photoshop Document), которые сохраняют всю структуру многослойного проекта для последующего редактирования.

\newpage
\refstepcounter{section}
\addcontentsline{toc}{section}{\protect\numberline{\thesection}Основы редактирования графических данных в прикладных программах}
\begin{center}\textbf{\thesection\ ОСНОВЫ РЕДАКТИРОВАНИЯ ГРАФИЧЕСКИХ ДАННЫХ В ПРИКЛАДНЫХ ПРОГРАММАХ}\end{center}
\vspace*{-1.5em}
\ManualSubsection{Принципы и среды редактирования графических данных}

Редактирование графических данных основывается на двух фундаментально различных подходах: разрушающем и неразрушающем. Разрушающее редактирование, характерное для ранних редакторов, предполагает прямое изменение пикселей исходного изображения, что делает процесс необратимым без наличия резервной копии. В противоположность этому, современные графические системы, такие как Adobe Photoshop и GIMP, делают акцент на неразрушающем редактировании \cite{adobe-helpx, gimp-docs}. Этот подход реализуется через использование слоёв, корректирующих масок и смарт-объектов, которые позволяют применять трансформации и эффекты, сохраняя исходные данные в неизменном виде. Такая методология обеспечивает высокую гибкость, возможность нелинейного монтажа и экспериментов без риска утраты качества, что стало стандартом в профессиональной обработке изображений.

\ManualSubsection{Базовые операции тоновой и цветовой коррекции}

Тоновая коррекция является первоочередным этапом обработки изображения и направлена на исправление распределения яркости. Ключевыми операциями выступают регулировка яркости и контраста. Яркость определяется интенсивностью пикселя, а её математическая коррекция может быть представлена как добавление постоянного значения $\Delta B$ к каждому пикселю $I(x, y)$. Коррекция контраста, влияющая на воспринимаемую разницу между светлыми и тёмными областями, обычно реализуется через масштабирование значений интенсивности вокруг среднего. Критическим инструментом анализа служит гистограмма распределения яркостей, визуализирующая тоновый диапазон снимка и позволяющая диагностировать проблемы, такие как недодержка или передержка \cite{rucomp}.

Цветокоррекция включает настройку баланса белого для компенсации цветовой температуры источника света и регулировку насыщенности. Баланс белого математически сводится к индивидуальному масштабированию каналов R, G и B для достижения нейтрального оттенка в заведомо белых участках сцены. Работа с насыщенностью, часто выполняемая в перцепционных цветовых пространствах вроде HSV, позволяет усиливать или ослаблять чистоту цвета, воздействуя на параметр S (Saturation). Эти операции являются основой для приведения изображения к естественному и гармоничному виду \cite{adobe-helpx}.

\ManualSubsection{Геометрические преобразования изображений}

Геометрические преобразования позволяют изменять пространственную композицию и ориентацию изображения. К базовым операциям относятся масштабирование, поворот, кадрирование и отражение. Масштабирование, описываемое умножением координат на коэффициенты $s_x$ и $s_y$, требует применения алгоритмов интерполяции для вычисления значений пикселей в новых координатах; наиболее качественные результаты обеспечивает бикубическая интерполяция. Поворот изображения на угол $\theta$ реализуется с помощью аффинного преобразования координат относительно центра вращения, что часто приводит к необходимости увеличения размеров холста для сохранения всех данных. Кадрирование представляет собой операцию выделения прямоугольной области, а отражение — зеркальное преобразование относительно вертикальной или горизонтальной оси. Современные редакторы выполняют эти преобразования неразрушающим способом, используя смарт-объекты или корректирующие слои-маски \cite{rucomp}.

\ManualSubsection{Основы работы в растровых редакторах: интерфейс и базовые техники}

Интерфейс современных растровых редакторов, таких как Adobe Photoshop, строится вокруг нескольких ключевых компонентов: рабочей области (холста), панели инструментов, палитры слоёв и панели настроек \cite{adobe-helpx}. Концепция слоёв является центральной, позволяя организовывать элементы композиции на независимых плоскостях, управлять их прозрачностью, режимами наложения и видимостью. К базовым техникам работы относится выделение областей с помощью инструментов прямоугольного и эллиптического выделения, лассо, а также «Волшебной палочки», которая автоматически выделяет области схожего цвета. Для неразрушающего скрытия частей изображения используются слоевые маски, работающие по принципу «чёрное — скрывает, белое — показывает». Инструменты ретуши, такие как «Штамп» и «Восстанавливающая кисть», применяются для устранения дефектов путём клонирования текстур с соседних участков \cite{gimp-docs}.

\ManualSubsection{Основы работы в векторных редакторах: принципы и инструменты}

Векторные редакторы, например Adobe Illustrator, оперируют не пикселями, а математическими описаниями контуров на основе кривых Безье. Основным инструментом для их создания и редактирования является «Перо», позволяющее размещать опорные точки и управлять направляющими ручками для точного контроля формы кривой \cite{adobe-helpx}. Для построения простых форм используются инструменты примитивов (прямоугольник, эллипс, многоугольник). Важным аспектом работы является редактирование узлов контура с помощью инструмента «Прямое выделение», которое позволяет перемещать точки, изменять тип узла (сглаженный, угловой) и регулировать кривизну сегментов. Объекты в векторной графике обладают атрибутами заливки и обводки, которым можно назначать сплошные цвета, градиенты или узоры. Логические операции над контурами (объединение, пересечение, вычитание) позволяют создавать сложные формы из простых составляющих.

\ManualSubsection{Композиция, работа с текстом и подготовка к публикации}

Создание комплексных графических композиций в любой среде основывается на грамотной организации слоёв, их группировке и использовании режимов наложения для взаимодействия визуальных элементов. Работа с текстом предполагает не только его ввод и форматирование, но и преобразование в кривые (контуры) для обеспечения независимости от наличия шрифтов при передаче файла, что критически важно в полиграфии и веб-дизайне \cite{adobe-helpx}. Финальным этапом является экспорт, где выбор формата файла определяется целевой платформой: векторные форматы (SVG, PDF) используются для масштабируемой графики и печати, растровые (JPEG, PNG, WebP) — для публикации в цифровой среде. При экспорте необходимо учитывать разрешение, цветовой профиль и параметры сжатия для оптимального баланса между качеством и размером итогового файла.

\newpage
\refstepcounter{section}
\addcontentsline{toc}{section}{\protect\numberline{\thesection}Цифровая обработка изображений и улучшение их качества}
\begin{center}\textbf{\thesection\ ЦИФРОВАЯ ОБРАБОТКА ИЗОБРАЖЕНИЙ И УЛУЧШЕНИЕ ИХ КАЧЕСТВА}\end{center}
\vspace*{-1.5em}
\ManualSubsection{Основные понятия и этапы цифровой обработки изображений}

Цифровая обработка изображений представляет собой научно-техническую дисциплину, занимающуюся анализом и модификацией визуальных данных с помощью алгоритмических и вычислительных методов. Фундаментальным объектом обработки является цифровое изображение, которое можно представить в виде двумерного дискретного массива пикселей, каждый из которых характеризуется координатами и значением интенсивности. Для цветных изображений чаще всего используется аддитивная модель RGB. Ключевыми характеристиками, определяющими техническое качество, являются разрешение (количество пикселей по осям) и глубина цвета (битность на канал), которые задают пределы детализации и цветового охвата \cite{gonzalez2018digital}. Типичный конвейер обработки включает последовательные этапы: захват изображения сенсором, предварительную обработку для коррекции базовых искажений и улучшения качества, сегментацию на семантически значимые области, анализ для извлечения информации и финальную визуализацию результатов.

\ManualSubsection{Классические методы улучшения контраста и тональной коррекции}

Одной из первостепенных задач улучшения качества является коррекция тонального диапазона и контраста изображения. Наиболее эффективным инструментом для анализа распределения яркости служит гистограмма — график, отображающий частоту встречаемости каждого уровня интенсивности. На основе её анализа применяются различные методы тональной коррекции. Эквализация гистограммы представляет собой нелинейное преобразование, направленное на перераспределение значений интенсивности для достижения равномерного (или заданного) распределения, что позволяет расширить динамический диапазон и улучшить восприятие деталей в тенях и светах. Линейное контрастирование является более простым методом, заключающимся в применении аффинного преобразования ко всем пикселям для растяжения или сжатия диапазона яркостей. Гамма-коррекция, как нелинейная операция, используется для компенсации нелинейности устройств отображения и позволяет независимо корректировать видимость деталей в тёмных и светлых участках сцены \cite{opencv2023tutorials}.

\ManualSubsection{Методы подавления шума и повышения резкости}

Шум, проявляющийся в виде случайных вариаций яркости или цвета, является частым артефактом, ухудшающим качество изображения. Для его подавления используются как линейные, так и нелинейные фильтры. К линейным относится усредняющий (боксовый) фильтр и фильтр Гаусса, которые реализуют свёртку изображения с соответствующим ядром, что эффективно для гауссовского шума, но приводит к размытию границ. Более совершенным нелинейным методом является медианная фильтрация, демонстрирующая высокую эффективность против импульсного шума типа «соль-и-перец» при лучшем сохранении резких переходов. Задача повышения резкости противоположна задаче шумоподавления и направлена на подчёркивание контуров и текстур. Классические методы основаны на использовании операторов, чувствительных к перепадам яркости, таких как фильтр Собеля для выделения границ или оператор Лапласа. Техника нерезкого маскирования (Unsharp Masking) является стандартным подходом, при котором к исходному изображению добавляется его высокочастотная компонента, полученная вычитанием размытой копии \cite{gonzalez2018digital, opencv2023tutorials}.

\ManualSubsection{Современные подходы на основе искусственного интеллекта}

Развитие глубокого обучения, в частности свёрточных нейронных сетей (CNN), произвело революцию в области улучшения качества изображений, предоставив методы, превосходящие классические алгоритмы. Нейросетевые модели успешно решают задачи сверхразрешения (Single Image Super-Resolution — SISR), позволяя восстанавливать высокочастотные детали при увеличении размера изображения, а также задачи обесшумливания (Denoising) и удаления артефактов сжатия (например, JPEG-артефактов). Генеративно-состязательные сети (GAN) открыли новые возможности для photo-realistic улучшений, таких как генерация реалистичных текстур при сильном апскейлинге, автоматическое раскрашивание чёрно-белых снимков и даже восстановление сильно повреждённых или фрагментированных изображений. Эти модели обучаются на больших наборах пар изображений «низкого» и «высокого» качества, учась восстанавливать семантически корректную информацию, что делает их инструментом не только для пиксельной, но и для семантической реставрации \cite{learnopencv2023ai}.

\ManualSubsection{Практическое применение в различных областях}

Методы цифровой обработки и улучшения изображений находят широкое и критически важное применение в разнообразных профессиональных сферах. В медицине они используются для повышения диагностической ценности рентгенограмм, снимков компьютерной и магнитно-резонансной томографии, позволяя визуализировать малоконтрастные структуры. В астрономии алгоритмы применяются для обработки данных телескопов, подавления атмосферных искажений и повышения разрешения снимков небесных тел. В криминалистике и системах видеонаблюдения технологии улучшения позволяют распознавать лица и номера на зашумлённых или низкокачественных записях. В области культурного наследия и фотографии методы используются для цифровой реставрации и колоризации исторических снимков. В робототехнике и системах компьютерного зрения предобработка изображений является обязательным этапом для обеспечения корректной работы алгоритмов навигации, обнаружения объектов и сценарного анализа \cite{gonzalez2018digital}.



\newpage
\refstepcounter{section}
\addcontentsline{toc}{section}{\protect\numberline{\thesection}Организация и хранение графических коллекций в цифровой среде}
\begin{center}\textbf{\thesection\ ОРГАНИЗАЦИЯ И ХРАНЕНИЕ ГРАФИЧЕСКИХ КОЛЛЕКЦИЙ В ЦИФРОВОЙ СРЕДЕ}\end{center}
\vspace*{-1.5em}
\ManualSubsection{Проблематика цифрового хаоса и принципы организации}

Экспоненциальный рост объёма графических данных в процессе профессиональной или учебной деятельности приводит к состоянию «цифрового хаоса», характеризующегося значительными временными затратами на поиск файлов, стрессом и риском безвозвратной утраты цифровых активов из-за сбоев оборудования или человеческой ошибки. Эфемерность цифровых объектов, их зависимость от конкретной программно-аппаратной среды и подверженность несанкционированному изменению делают вопросы организации и сохранности не второстепенными, а базовыми профессиональными компетенциями. Целью построения системы управления коллекцией является обеспечение возможности быстрого поиска (в идеале в течение 30 секунд) любого материала и гарантия его долговременной сохранности, что подразумевает не только физическую целостность данных, но и сохранение их аутентичности, читаемости и полезности в будущем \cite{oais_reference_model}.

Эффективная система строится на двух взаимодополняющих принципах: структуре папок («скелет») и работе с метаданными («душа»). Иерархическая структура папок создаёт понятный макроуровневый порядок и логическую группировку файлов, обеспечивая визуальную навигацию. Метаданные, встроенная в файлы описательная информация, выступают в роли электронного каталога, позволяя осуществлять семантический поиск по содержанию (кто, что, где, когда) без знания физического расположения файла. Папки без метаданных превращаются в кучи файлов для ручного перебора, а метаданные без структуры — в виртуальную свалку, поэтому оба принципа должны применяться в тандеме \cite{principles_archival_description}.

\ManualSubsection{Стратегии организации файловой структуры}

Выбор логики построения иерархии папок определяется спецификой деятельности и характером графической коллекции. Хронологический принцип, при котором структура выстраивается по пути «Год → Месяц → Событие», является простым, надёжным и исключающим дублирование, что делает его идеальным для фотографов и личных архивов. Проектный принцип, организующий файлы в папки типа «01\_Бриф», «02\_Референсы», «03\_Рабочие\_файлы», позволяет сосредоточить всю информацию, относящуюся к конкретной задаче, в одном месте, что оптимально для дизайнеров и студентов. Тематический принцип, например, «Референсы → UI\_UX → Мобильные\_приложения», хорошо подходит для коллекционеров визуальных материалов, хотя и сталкивается с проблемой отнесения одного файла к нескольким темам одновременно. Независимо от выбранной стратегии, архивная практика рекомендует документировать организационную структуру для обеспечения её понятности в долгосрочной перспективе \cite{principles_archival_description}.

\ManualSubsection{Метаданные как основа для поиска и долговременной сохранности}

Метаданные графических файлов подразделяются на автоматически записываемые и добавляемые вручную. К первым относится информация в формате EXIF (Exchangeable Image File Format), содержащая технические параметры съёмки (дата, время, параметры камеры, GPS-координаты), полезные для анализа, но малопригодные для семантического поиска. Основой для «умного» поиска выступают метаданные IPTC (International Press Telecommunications Council), которые пользователь добавляет самостоятельно: ключевые слова (теги), заголовок, описание и информация об авторе \cite{iipa_metadata_handbook}.

Правильное заполнение метаданных, в частности тегов, должно давать ответы на вопросы «Кто? Что? Где? Когда?» и включать как конкретные термины, так и концептуальные или эмоциональные описания. Для обеспечения долговременной сохранности и установления контекста критически важно фиксировать такие метаданные, как полное имя автора, дата создания, название проекта и информация о лицензии, непосредственно в файл, а не только в базу данных каталогизатора, чтобы избежать их потери при сбое программного обеспечения \cite{oais_reference_model, iipa_metadata_handbook}.

\ManualSubsection{Программные инструменты для каталогизации и управления}

Управление крупными графическими коллекциями через стандартный файловый менеджер неэффективно, что обуславливает необходимость использования специализированного программного обеспечения. На начальном уровне достаточно мощным решением является Adobe Bridge, предоставляющий инструменты для просмотра, пакетного редактирования и добавления метаданных. Для профессиональной работы с фотографиями «золотым стандартом» выступает Adobe Lightroom Classic, который представляет собой интегрированную среду для импорта, каталогизации, неразрушающего редактирования и экспорта. Его работа основана на каталоге — базе данных, хранящей все метаданные, рейтинги и историю редактирования, в то время как исходные файлы остаются на своих местах в файловой системе \cite{iipa_metadata_handbook}. Важным правилом является управление файлами (перемещение, переименование) только через интерфейс самой программы-каталогизатора во избежание разрыва связей между каталогом и физическими файлами, что в архивной практике приравнивается к нарушению целостности коллекции \cite{oais_reference_model}.

\ManualSubsection{Стратегия надёжного хранения и резервного копирования}

Поскольку любые носители данных имеют ограниченный срок службы и подвержены рискам физического повреждения, кражи или стихийных бедствий, единственной гарантией сохранности цифровой коллекции является продуманная стратегия резервного копирования. Эмпирическое правило «3-2-1» является отраслевым стандартом: необходимо иметь три полные копии данных, хранящиеся на двух разных типах носителей, причём одна из копий должна располагаться географически удалённо (оффсайт). Например, рабочая копия может находиться на SSD ноутбука, первая резервная — на внешнем жёстком диске дома, а вторая — в облачном хранилище. Простого копирования недостаточно; для долгосрочного архивирования требуется периодическая миграция данных на новые носители и верификация их целостности с помощью контрольных сумм (хеш-сумм) \cite{oais_reference_model}.

\ManualSubsection{Практические рекомендации по именованию и версионности}

Согласованное именование файлов является простым, но мощным организационным инструментом. Рекомендуется использовать информативные имена по шаблону «ГГГГММДД\_Проект\_Описание\_Версия.расширение» (например, «20240521\_Курсовой\_Лендинг\_v3.psd»), что обеспечивает автоматическую хронологическую сортировку и мгновенное понимание содержания файла. Для управления процессом работы важна практика версионности: сохранение ключевых промежуточных версий файлов с инкрементной нумерацией (v1, v2, final) вместо перезаписи, что позволяет отслеживать изменения и возвращаться к предыдущим состояниям. Для фиксации изменений на уровне данных, особенно при долгосрочном хранении, рекомендуется использование контрольных сумм. Регулярность поддержания порядка — выделение коротких промежутков времени (например, 15 минут в неделю) на организацию новых файлов, добавление метаданных и выполнение резервных копий — является залогом устойчивости всей системы в долгосрочной перспективе \cite{principles_archival_description}.

\newpage
\section*{}
\vspace*{-2.5em}
\begin{center}\textbf{ЗАКЛЮЧЕНИЕ}\end{center}
\phantomsection
\addcontentsline{toc}{section}{Заключение}

Технологии создания и обработки графических данных представляют собой комплексную дисциплину, лежащую в основе современной визуальной коммуникации, цифрового искусства, дизайна и научной визуализации. В ходе проведённого исследования были проанализированы фундаментальные принципы представления графической информации, её обработки, хранения и управления.

Установлено, что базовым выбором для любой графической задачи является определение математической модели изображения. Растровая графика, основанная на пиксельной матрице, является незаменимой для передачи фотореалистичных изображений и сложных тоновых переходов, в то время как векторная графика, оперирующая математическими описаниями контуров, обеспечивает бесконечную масштабируемость и точность, что критически важно для полиграфии, веб-дизайна и создания фирменного стиля.

Ключевым аспектом эффективной работы с графикой является понимание и правильное применение цветовых моделей. Аддитивная модель RGB является стандартом для экранного представления, субтрактивная CMYK — для полиграфического производства, а перцепционные модели (HSL, Lab) служат основой для интуитивной цветокоррекции и профессиональной ретуши. Выбор формата файла и метода сжатия (с потерями или без потерь) представляет собой поиск оптимального баланса между качеством визуальных данных, размером файла и требованиями целевой платформы, будь то веб-публикация, печать или долгосрочное архивное хранение.

Процесс редактирования графических данных, будь то в растровых или векторных редакторах, базируется на современных принципах неразрушающего редактирования через использование слоёв, масок и корректирующих эффектов, что обеспечивает гибкость и сохранность исходного материала. Задачи цифровой обработки изображений, от тоновой коррекции и подавления шума до сложного восстановления с помощью алгоритмов искусственного интеллекта, направлены на улучшение визуального качества и пригодности изображений для последующего анализа.

Финальным, но не менее важным звеном технологического цикла является организация и хранение графических коллекций. Построение логической файловой структуры, систематическое заполнение метаданных (IPTC, EXIF) и реализация надёжной стратегии резервного копирования по правилу «3-2-1» трансформируют хаотичное скопление файлов в профессионально управляемый цифровой актив, обеспечивая его долговременную сохранность, целостность и лёгкость поиска.

Таким образом, владение полным спектром технологий работы с графическими данными — от их генезиса до архивации — является неотъемлемой составляющей цифровой грамотности и профессиональной компетенции в любой визуально-ориентированной сфере, обеспечивая как эффективность текущей работы, так и сохранение цифрового наследия для будущего.

\newpage
\begin{center}\textbf{СПИСОК СОКРАЩЕНИЙ}\end{center}
\phantomsection
\addcontentsline{toc}{section}{Список сокращений}

\vspace{1em}

{
\small
\noindent
\begin{tabular}{|l|p{0.75\textwidth}|}
\hline
\textbf{Сокращение} & \multicolumn{1}{c|}{\textbf{Расшифровка}} \\
\noalign{\hrule height 1.2pt}
\hline
DPI & Количество точек на дюйм. Единица измерения разрешения печатных и отображаемых изображений.\\
\hline
RGB & Аддитивная цветовая модель, основанная на смешении красного, зелёного и синего цветов. Стандарт для экранного отображения. \\
\hline
CMYK & Субтрактивная цветовая модель, использующая голубой, пурпурный, жёлтый и чёрный цвета. Стандарт для полиграфической печати.\\
\hline
LAB (CIELAB) & Аппаратно-независимая перцепционная цветовая модель, разделяющая информацию о светлоте (L) и цветности (a, b). \\
\hline
EXIF & Формат метаданных, автоматически записываемых цифровыми камерами и содержащих технические параметры съёмки.\\
\hline
IPTC & Международный стандарт метаданных, предназначенный для ручного добавления описательной информации (ключевые слова, автор, описание) к изображениям.\\
\hline
SVG & Масштабируемая векторная графика. Открытый XML-формат для описания двумерной векторной и смешанной графики.\\
\hline
JPEG & Распространённый формат сжатия растровых изображений с потерями, оптимальный для фотографий.\\
\hline
PNG & Формат растровой графики, поддерживающий сжатие без потерь и альфа-канал (прозрачность). \\
\hline
RAW & «Сырой» формат данных, содержащий необработанную информацию с матрицы цифровой камеры.\\
\hline
PSD & Родной формат файлов Adobe Photoshop, сохраняющий слои, маски и другие элементы редактирования.\\
\hline
UI & Пользовательский интерфейс. Средства, обеспечивающие взаимодействие пользователя с программой или устройством.\\
\hline
UX & Пользовательский опыт. Общее восприятие и эмоции пользователя, возникающие при взаимодействии с продуктом.\\
\hline
\end{tabular}
}

\clearpage
\nocite{*}
\printbibliography[heading=bibliography]

\end{document}
